\ifthenelse{\equal{\lang}{ko}}{%

우리의 평가는 자동화된 AI 동향 감지 시스템의 강점과 한계를 모두 부각시키는 동시에, 중요한 윤리적 고려사항과 연구적 함의를 드러낸다. 

\subsection{강점(Strengths)}
본 프레임워크는 기존 모니터링 접근법에 비해 명확한 장점을 보여준다. 
실시간 탐지와 LLM 기반 의미 추론을 결합함으로써, 시스템은 다양한 과제에서 높은 분류 정확도를 유지하면서도 수작업 분석가의 업무를 80\% 이상 줄인다. 
모듈형 아키텍처는 새로운 LLM과 컴플라이언스 데이터셋과의 원활한 통합을 가능하게 하여 장기적 확장성을 보장한다. 
출처 메타데이터가 포함된 구조화된 출력은 해석 가능성을 개선하고 기업 의사결정에 대한 신뢰를 촉진한다. 
적용상의 이점뿐만 아니라, \textbf{규제 감지를 위한 LLM 기반 분류체계} 도입은 AI 거버넌스 정보학이라는 신흥 분야에 개념적 기여를 한다.

\subsection{한계(Limitations)}
그럼에도 불구하고 몇 가지 한계가 존재한다. 
첫째, 시스템 성능은 관할권별로 크게 차이가 나는 온라인 데이터 소스의 품질과 범위에 크게 의존한다. 
둘째, 다국어 검증을 포함했으나 저자원 언어에 대한 전면적 평가는 여전히 해결되지 않은 과제이다. 
셋째, 사용자 연구에서 고무적인 결과를 보였지만 참여자 집단(18명의 분석가)은 전 세계 규제 실무의 다양성을 충분히 반영하지 못할 수 있다. 
마지막으로, 동적 규제 환경과 진화하는 AI 거버넌스 프레임워크에서 장기적 안정성은 추가 검증이 필요하다.

\subsection{타당성 위협(Threats to Validity)}
우리는 세 가지 주요 타당성 위협을 식별하였다. 
(1) \textbf{데이터셋 편향:} 평가 코퍼스가 고자원 관할권(EU, 미국 등)을 과대표집했을 수 있어 일반화 가능성이 제한될 수 있다. 
(2) \textbf{사용자 연구 범위:} 분석가 참여자는 소수 기관에서 모집되어 사용성에 대한 인식이 왜곡될 수 있다. 
(3) \textbf{개념 드리프트:} 빠르게 변화하는 AI 규제가 분포적 변화를 초래하여 장기적 정확도를 낮출 수 있다. 
이를 완화하기 위해서는 지속적인 재학습과 데이터셋 업데이트가 필요하다. 

\subsection{윤리적 및 사회적 고려사항}
거버넌스를 위한 자동화된 감지는 윤리적 위험을 수반한다. 
학습 데이터의 편향이 규제 해석으로 전파되어 민감한 조항이 잘못 분류되거나 소외된 관할권이 과소 대표될 위험이 있다. 
출력이 컴플라이언스 결정에 활용될 경우, 투명성과 설명 가능성은 이해관계자의 신뢰를 보장하는 데 필수적이다. 
또한 대규모 모니터링은 적절한 보호장치 없이는 감시(surveillance)와 유사하게 보일 수 있어 프라이버시 위험도 관리해야 한다. 
본 설계는 구조화된 출력, 편향 감사, analyst-in-the-loop 재학습을 통해 이러한 위험을 완화하지만, 책임 있는 배포를 위한 보다 광범위한 표준이 필요하다. 

\subsection{산업적 사례 스냅샷 (익명화)}
실제 영향을 평가하기 위해, 기밀을 유지하면서 두 개의 파일럿 배포를 수행하였다.\footnote{조직명은 NDA에 따라 비공개 처리됨.}
\textbf{사례 A (금융, EU)}: 내부 규제 공지 수집과 보고가 자동화되었다. 평균 엔드투엔드 지연은 3.5초 이하로 유지되었고, 업무량 감소는 평가 결과(≈80\%+)와 유사하였다. 분석가들은 누락 항목이 줄고 에스컬레이션 속도가 빨라졌다고 보고하였다.  
\textbf{사례 B (공공기관, 한국)}: 다국어 포털의 주간 컴플라이언스 브리핑이 증거 기반 인용을 포함한 구조화된 요약으로 통합되었다. 보고 준비 시간이 약 70\% 단축되었고, 출처 로그를 통해 교차 검증 가능한 커버리지 일관성이 개선되었다.  
이러한 사례들은 통합된 탐지–추론–보고 스택이 기존 거버넌스 워크플로우에 최소한의 적응만으로도 이식 가능함을 입증한다.

\subsection{시사점(Implications)}
본 연구 결과는 학계와 산업 모두에 시사점을 제공한다. 
연구자에게는 규제 감지를 대규모 언어 모델을 위한 새로운 벤치마크 과제로 확립하여, 다국어, 공정성 인식, 설명 가능한 접근법에 대한 심층 탐구를 장려한다. 
실무자에게는 기업 컴플라이언스를 위한 확장 가능한 인프라를 제공하여, 조직이 실시간으로 규제 동향을 추적할 수 있게 한다. 
기회와 위험을 모두 조명함으로써, 본 연구는 자연어 처리, 법학, 책임 있는 AI의 교차점에서 \textbf{AI 거버넌스 정보학}이라는 새로운 연구 영역의 형성에 기여한다. 


}{%
% -------------------- English Version --------------------


Our evaluation highlights both the strengths and limitations of the Automated AI Trend Sensing System, while also surfacing important ethical considerations and research implications. 

\subsection{Strengths}
The framework demonstrates clear advantages over traditional monitoring approaches. 
By combining real-time detection with LLM-driven semantic reasoning, the system reduces manual analyst effort by over 80\% while maintaining high classification accuracy across diverse tasks. 
Its modular architecture allows seamless integration with emerging LLMs and compliance datasets, ensuring long-term scalability. 
Structured outputs, accompanied by provenance metadata, further improve interpretability and foster trust in enterprise decision-making. 
Beyond applied benefits, the introduction of an \textbf{LLM-driven taxonomy for regulatory sensing} represents a conceptual contribution to the emerging field of AI governance informatics.

\subsection{Limitations}
Despite these contributions, several limitations remain. 
First, system performance depends heavily on the quality and coverage of online data sources, which vary significantly across jurisdictions. 
Second, while cross-lingual validation was included, full-scale evaluation across low-resource languages remains an open challenge. 
Third, although user studies showed promising results, the participant pool (18 analysts) may not capture the diversity of perspectives in global regulatory practice. 
Finally, long-term stability under dynamic regulatory environments and evolving AI governance frameworks must be further validated.

\subsection{Threats to Validity}
We identify three primary threats to validity. 
(1) \textbf{Dataset bias:} The evaluation corpus may overrepresent high-resource jurisdictions (e.g., EU, U.S.), potentially limiting generalizability. 
(2) \textbf{User study scope:} Analyst participants were recruited from a small number of organizations, which may skew perceptions of usability. 
(3) \textbf{Concept drift:} Rapidly changing AI regulations introduce distributional shifts that may reduce long-term accuracy. 
Mitigation requires continuous retraining and dataset updates. 

\subsection{Ethical and Societal Considerations}
Automated sensing for governance introduces ethical risks. 
Biases in training data may propagate to regulatory interpretations, risking misclassification of sensitive clauses or underrepresentation of marginalized jurisdictions. 
Transparency and explainability are critical to ensuring stakeholder trust, particularly when outputs inform compliance decisions. 
Privacy risks must also be managed: large-scale monitoring could resemble surveillance without safeguards. 
Our design mitigates these risks through structured outputs, bias auditing, and analyst-in-the-loop retraining, but broader standards for responsible deployment remain necessary. 

\subsection{Industrial Case Snapshots (Anonymized)}
We conducted two pilot deployments to assess real-world impact while preserving confidentiality.\footnote{Organization names withheld under NDA.}
~\\
\textbf{Case A (Finance, EU)}: Internal regulatory bulletin ingestion and reporting were automated. Median end-to-end latency remained sub-3.5s with workload reduction comparable to our evaluation (≈80\%+). Analysts reported fewer missed items and faster escalations, aligning with our measured efficiency gains:contentReference[oaicite:1]{index=1}.
~\\
\textbf{Case B (Public Agency, KR)}: Weekly compliance briefs (multi-lingual portals) were consolidated into structured summaries with evidence-linked citations. Report preparation time decreased by $\sim$70\% and coverage consistency improved (cross-checkable through provenance logs), mirroring our main results:contentReference[oaicite:2]{index=2}.
These snapshots corroborate industrial viability: the unified sensing--reasoning--reporting stack transfers with minimal adaptation to existing governance workflows.


\subsection{Implications}
Our findings have implications for both research and practice. 
For researchers, the results establish regulatory sensing as a new benchmark task for large language models, encouraging deeper exploration of cross-lingual, fairness-aware, and explainable approaches. 
For practitioners, the system provides an extensible infrastructure for enterprise compliance, enabling organizations to track regulatory developments in real time and at scale. 
By highlighting both opportunities and risks, this work contributes to the formation of \textbf{AI governance informatics}—a new domain at the intersection of NLP, law, and responsible AI. 


}
